\section{Deadline Miss Detection}
In our system, each task is associated with an hard deadline, as described in
Table \ref{tab:tasks-list}. In order to detect a deadline miss, since the 
\textit{RTIC} framework does not provide any built-in mechanism for this, for
each task we implemented a watchdog task responsible for monitoring its
deadline. \\ 
To achieve the desired behaviour, i.e., a watchdog declares a deadline miss if
and only if the associated task fails to complete its execution before its
deadline, we used a shared resource, called
\texttt{deadline\_protected\_object}. This resource is a struct containing a
boolean flag, which value rappresents whether the task has completed its
execution or not. In \textit{RTIC} framework, each shared resource needs
to be acquired before being accessed or before invoking any of its associated
methods. Therefore, when a monitored task acquired the lock on the
\texttt{deadline\_protected\_object} to invoke \texttt{cancel\_deadline}, the
method responsible for acknowledging the deadline, it can not be preempeted by
its watchdog task. This prevents the watchdog from preempting the monitored task
while it is updating the flag, avoiding undesired deadline miss detections. \\
Additionaly, the \texttt{deadline\_protected\_object} struct contains an
activations counter, which allows for finer control over the requests to clear
deadline misses, allowing stale requests to be ignored. \\
Since in our system we have periodic and sporadic tasks, we implemented two
different types of watchdogs: 
\begin{itemize}
    \item \texttt{periodic\_deadline\_watchdog}: this watchdog has as locals are
    the period and the next deadline of the monitored task. It delays itself
    until the next deadline and then, after acquiring the shared
    \texttt{deadline\_protected\_object}, it invokes on it the method
    \texttt{deadline\_miss\_detected}. This method checks the flag and, if it is
    still false, it declares a deadline miss. Then, the watchdog increments the
    next deadline by the period of the monitored task and it loops again;
    \item \texttt{sporadic\_deadline\_watchdog}: this watchdog waits on an
    activation signal that notices it of a new activation of the monitored task. 
    The watchdog then calculates the next deadline by adding the relative
    deadline with the activation time of the monitored task, passed as argument
    inside the precedent signal. The watchdog then delays itself until the
    obtained absolute deadline and, after acquiring the shared
    \texttt{deadline\_protected\_object}, it invokes on it the method 
    \texttt{deadline\_miss\_detected}, as happens in the \\
    \texttt{periodic\_deadline\_watchdog} case. After that, it loops again,
    waiting for the next activation signal.
\end{itemize}
To avoid the first activation of a periodic watchdog from starting before the
monitored task can actually execute due to system initialisation, 
we set its first deadline equal to the system activation instant plus the
task's relative deadline. 
The system activation instant consists in a static \texttt{AtomicU32}
rappresenting, in ticks, the instant befoore which all initializtion-spawned
tasks are delayed: this mechanisms provides the system with a well defined
"\textit{zero time}". \\ 
Moreover, in order to provide accurate deadline management
for sporadic tasks, their activation signals occur inside a critical section.
This prevents unwanted preemptions between the task activation and the sending
of the signal to the corresponding watchdog, which could otherwise lead to an
incorrect calculation of the next deadline.  